\section{Scope, Method, and Qualifications}

\subsection{Reader, question, and claim boundary}

This is a close exegetical study of one passage---John 20:3--10, with Luke
24:12---for a reader who wants to know what the text says and what can honestly
be built on it. No Greek, Latin, or Hebrew is presupposed; every word discussed
is transliterated, glossed from a named lexicon, and traced through its other
occurrences.

Its governing question is a claim, not a thesis: that St.\ Peter knew Jesus had
risen because he noticed how the burial cloths were left. The study tests that
claim, and reaches a divided verdict, stated in full in section~9---the
observation is Peter's on the Fourth Gospel's own showing; the inference is
attributed by John to the other disciple and by Luke to nobody.

Two subsidiary questions follow from it and are worked to the same standard:
whether the ancient argument that the arrangement tells against grave-robbery
is sound, and whether the modern claim that a folded napkin was a Jewish
master-and-servant signal has any attestation. The second is answered in the
negative, with the search described.

The study does not argue for or against the Resurrection, does not assess any
relic, and does not adjudicate the historicity of the empty tomb. Section~10
states those boundaries and the four questions of fact left open.

\subsection{Terminology fixed for this study}

\begin{itemize}
\item \textbf{The claim under test} means the sentence quoted at the head of
section~1, in the maintainer's own words. It is treated as an object of
examination throughout and is nowhere assumed.
\item \textbf{The other disciple} is used throughout for the figure John calls
\emph{ho allos math\=et\=es}. The traditional identification with John son of
Zebedee is not required by any argument here and is not assumed; where a source
makes it, the source's language is preserved.
\item \textbf{Gloss} marks an explanation of a word's range from a named lexical
authority or from its occurrences. It is never a translation of Scripture. No
translation of Scripture has been composed for this study.
\item \textbf{Working translation} marks an English rendering of a Latin or
Greek clause supplied by this study where no public-domain published
translation of that clause was consulted. Every such rendering is labelled at
its own place; none is offered for recitation or as a version of Scripture.
\item \textbf{Transliteration} is mechanical and is not a translation. Greek
long vowels are marked \emph{\=e} and \emph{\=o}; chi is \emph{ch}; upsilon is
\emph{y} outside diphthongs. Hebrew and Aramaic are transliterated without
scholarly diacritics because this document carries no Hebrew face.
\item \textbf{Material negative result} means a case in which a witness that had
the occasion to make an argument, and would be expected to make it, does not.
Such results are reported as findings.
\end{itemize}

\subsection{Source hierarchy and evidence classes}

The governing texts are the Greek New Testament, the Latin Vulgate, and the
English Douay--Rheims, in that order of authority for the wording and in the
reverse order of accessibility for the reader.

\textbf{Greek.} Robinson and Pierpont's Byzantine Textform (2018) in the
registered electronic edition of the byztxt CSV files. This is one textual
tradition, not a critical eclectic text. Where its apparatus records a
Nestle--Aland divergence at a verse used here, the divergence is stated. No
Nestle--Aland or United Bible Societies edition was collated; where the modern
critical text is mentioned, the statement rests on what the registered
apparatus records and is labelled as such.

\textbf{Latin.} The Clementine Vulgate in Michael Hetzenauer's 1914 Pustet
printing, read from the registered page images of SacredBible's photographed
copy. Five passages were read on the page and transcribed: Jn 20:3--10, Jn
19:39--40, Jn 11:44, Lk 23:53, Lk 24:12. Hetzenauer's editorial marginal
summaries, section headings, and cross-reference letters are that edition's
apparatus and are identified as such wherever quoted.

\textbf{English Scripture.} The public-domain Douay--Rheims in the Challoner
revision, from the registered electronic edition, in Vulgate numbering
throughout. The Authorized Version is quoted twice, from the Project Gutenberg
text, solely to establish what English Bibles say at Jn 20:7; it carries no
authority in this study.

\textbf{Lexical.} Joseph Henry Thayer's revision of Grimm's \emph{Clavis Novi
Testamenti} (corrected edition, American Book Company), cited by headword and
page and read at those pages in a facsimile. It is a nineteenth-century lexicon
used as an identified working authority. No modern lexicon (Bauer--Danker--Arndt,
Liddell--Scott--Jones in a current edition, or the \emph{Diccionario
Griego-Espa\~nol}) was collated; where a delegated search reported their
readings, the readings were not relied on for any claim in the text.

\textbf{Patristic and medieval.} Read at their own loci, never through a later
author's quotation, with one labelled exception: the \emph{Catena aurea} is by
its nature a quotation collection and is identified as such. Chrysostom and
Augustine were read in page images of the printed Nicene and Post-Nicene
Fathers volumes. Augustine's Latin, Gregory's Latin, Bede's Latin, and Ambrose's
Latin were read in electronic transcriptions of Migne with the transcription
ceiling recorded at each place. Cyril was read in an electronic transcription of
the Library of Fathers translation. Theophylact was read in the Internet
Archive's OCR of Migne's facing Latin column, and is reported at an OCR ceiling;
his Greek was not collated by this study. Aquinas's \emph{Lectura} was read in
the Busa--Alarc\'on Latin at the Corpus Thomisticum.

\textbf{Rabbinic.} Searched and read through Sefaria's public text and search
interfaces. Quotations are taken only from versions whose recorded licence
permits it: the public-domain Torat Emet Hebrew, A.~Cohen's public-domain 1921
Cambridge \emph{Tractate Berakot}, and the CC0 Sefaria Community Translation.
The William Davidson English is CC BY-NC and is not quoted anywhere in this
study.

\textbf{Modern secondary.} Westcott's commentary, Latham's \emph{Risen Master},
Westcott and Hort's Appendix, Smith's \emph{Dictionary}, Lightfoot's
\emph{Horae}, and Suetonius in Rolfe's Loeb are all public-domain works read
directly. Present-day web pages---\emph{Jerusalem Perspective}, TruthOrFiction,
GotQuestions, and one 2007 mailing-list archive---are used as dated evidence of
what the modern claim says and when it circulated, quoted in short clauses with
attribution, and are under their own rights.

\subsection{Method, and the use of delegated research}

Research proceeded claim by claim. Each consequential statement in the text was
traced to a witness read for this study; where a witness could not be reached, a
ceiling is recorded rather than a substitute cited.

Part of the search work was delegated to assisting research processes, and this
should be visible rather than hidden. The findings so obtained were treated as
leads, and every finding relied on in the argument was independently re-checked
at its own source before use: the Mishnah and Talmud passages were re-retrieved
and the Sefaria phrase searches re-run for this study; the \emph{Jerusalem
Perspective} page, the TruthOrFiction and GotQuestions pages, the 2007
mailing-list archive, Smith's \emph{Mantele} article, Suetonius, Cyril,
Gregory, Ambrose, Bede on Luke, Theophylact's Latin column, Aquinas's
\emph{Lectura}, Jerome, and Lightfoot's four volumes were all fetched and read
directly. Three findings are reported at an explicitly delegated ceiling and are
marked at their own places: the absence of the stem \emph{sudari-} in Migne's
Bede homiliary, the reading of the \emph{Gospel of Peter} in M.~R. James's
translation, and the reported treatments in modern copyrighted Johannine
commentaries, which are not relied on for any claim.

Negative results were sought deliberately and are reported with their
boundaries. A phrase search cannot exclude a synonym, an unindexed work, a
paraphrase, or a language not searched, and the sections that report negatives
say so at the point of reporting.

\subsection{Material qualifications}

\begin{enumerate}
\item The judgment reached in section~9 is a judgment about what two
first-century texts attribute to a named man. It is not a judgment about
Peter's faith, his standing, or his subsequent testimony, all of which are
outside this study.
\item Chrysostom's and Cyril's collective reading of Jn 20:8 is a genuine
patristic position held by Fathers of the first rank; the study's disagreement
with it is exegetical and is argued, not asserted, and their position is stated
at full strength in the synthesis before it is answered.
\item Augustine's minimizing reading is not adopted outright. It is reported as
defensible, its two costs are stated, and the study does not decide between it
and the intermediate readings.
\item The transliteration of \emph{entetyligmenon} as ``wrapped'' or ``rolled''
rests on Thayer, on the word's other New Testament occurrences, and on the
Vulgate's choice of \emph{involutum}; Latham, a partisan of the arrangement
argument, is quoted making the same point against interest. The claim is that
``folded'' is not what the word says, not that no cloth in the tomb could have
been folded.
\item \emph{Eis hena topon} is left undecided among three readings; no argument
in the study depends on the choice.
\item Luke 24:12's textual standing is disputed and the dispute is recorded. No
conclusion rests on that verse alone.
\item The Latin quoted from Hetzenauer 1914 is that printing's text; other
Clementine printings may differ in accidentals, and this study collated none.
\item The rabbinic negative is bounded by the indexed corpus, the English
phrasings tried, and the date of the search (25 July 2026). It is stated as a
bounded and correctable negative in the synthesis block that reports it,
together with the strongest argument against it and the finding that would
overturn it.
\item The reconstruction of the modern claim's origins is bounded by what could
be reached. Google Books, HathiTrust, and the Usenet archives were unavailable
during this research; pre-2006 print circulation could not be checked. The
\emph{Jerusalem Perspective} publication date is that site's own metadata field,
and the earliest Internet Archive capture of the page postdates it by six years.
\item Quotations from present-day web pages and from David Bivin's response are
brief, attributed, and used as evidence of what was said and when; they are not
offered under this project's licence.
\item The English renderings of Latin from Ambrose, Gregory, Bede, Theophylact,
and Aquinas are this study's labelled working translations of the clauses
quoted. They are not offered as published translations and are not for
recitation. Where a public-domain published English translation exists and was
read---Chrysostom, Augustine, Cyril, Jerome---it is used instead and identified.
\item Every claim about what the tradition does \emph{not} say is a claim about
the witnesses named in section~7, at the loci named there. It is not a claim
about the whole of patristic literature.
\end{enumerate}

\subsection{Coverage, currentness, rights, and review}

No liturgical rite, calendar, jurisdiction, or mutable discipline governs any
conclusion here, and no canonical or legal claim is made. All online witnesses
were consulted on 25 July 2026; the two Vulgate page images newly registered for
this study were retrieved the same day, and the earlier page images of that
edition already in the repository were retrieved on 23 July 2026.

Scripture, patristic translations, third-party transcriptions, lexicon text, and
the rabbinic translations retain their own rights status and are not offered
under this project's licence. The public-domain basis of each quoted witness is
recorded with its source record. The Sefaria William Davidson English is
CC BY-NC and is quoted nowhere.

This is an internally source-audited working study produced by an AI-assisted
process. It has received no independent exegetical, patristic, rabbinic,
text-critical, or ecclesiastical review, and it claims no ecclesiastical
approval. It is a study aid, not a commentary of record.
